\documentclass[11pt]{article}
\usepackage{amsmath,amsthm,amssymb}
\usepackage[margin=1in]{geometry}
\usepackage{hyperref}
\usepackage{booktabs}
\usepackage{enumitem}

\newtheorem{theorem}{Theorem}
\newtheorem{lemma}[theorem]{Lemma}
\newtheorem{proposition}[theorem]{Proposition}
\newtheorem{corollary}[theorem]{Corollary}
\newtheorem{conjecture}[theorem]{Conjecture}
\theoremstyle{definition}
\newtheorem{definition}[theorem]{Definition}
\newtheorem{remark}[theorem]{Remark}
\newtheorem{observation}[theorem]{Observation}
\newtheorem{question}[theorem]{Question}

\title{No naive Hecke iso $B_{\ell+1}^{(\lambda)}V_\lambda \cong V_{\lambda^{(\ge 2)}}$:\\
  a precise computational refutation}
\author{Clio}
\date{13 May 2026 (evening)}

\begin{document}
\maketitle

\begin{abstract}
The May-13 closed-form theorem says
$\dim B_{\ell+1}^{(\lambda)} V_\lambda = f^{\lambda^{(\ge 2)}}$ in the
stable regime, where $\lambda^{(\ge 2)}$ is the decoration shape
(column~$1$ deleted). The natural next question is whether this is a
\emph{representation-theoretic} isomorphism $B_{\ell+1}^{(\lambda)}
V_\lambda \cong V_{\lambda^{(\ge 2)}}$ as $H_q(S_{n-\ell})$-modules
under some embedding of $H_q(S_{n-\ell})$ into $H_q(S_n)$.

\emph{We refute this in its most natural form.} For the stable test case
$\lambda = (3, 2, 1, 1, 1)$ at $n = 8$, where $\dim B = 2 = f^{(2,1)}$
and $\dim V_{\lambda^{(\ge 2)}} = \dim V_{(2,1)} = 2$, direct computation
in the Hoefsmit seminormal basis shows that the only Hecke generators
$T_j$ ($1 \le j \le n - 1$) that preserve $B$ are
\[
   T_1 \quad (\text{acting as scalar } -1)
   \qquad\text{and}\qquad
   T_\ell = T_5 \quad (\text{acting as scalar } q).
\]
All other $T_j$ \emph{fail to preserve} $B$ (they extend $B$ to a
strictly larger subspace of $V_\lambda$). Moreover, no product
$T_i T_j$ or braid triple $T_i T_j T_i$ preserves $B$ unless it is
already in $\langle T_1, T_5\rangle$. Hence no Hecke subalgebra
generated by simple $T_j$'s or their short products carries a
non-trivial $H_q(S_{n-\ell})$-action on $B$.

This forces one of three resolutions:
\begin{enumerate}[label=(\alph*),topsep=0.3em,itemsep=0.2em]
\item the iso, if it exists, is via \emph{exotic} operators in
  $H_q(S_n)$ (e.g.\ Murphy/Jucys elements or specific intertwiners) ---
  none of which were located in this session;
\item the iso holds only at the level of \emph{vector spaces} (the chain
  bijection $\mathrm{Chain}(\lambda) \leftrightarrow \mathrm{SYT}(\lambda^{(\ge 2)})$,
  established in the May-13 paper) and has no Hecke-module upgrade;
\item the iso holds via a \emph{subquotient} or \emph{induced} Hecke
  structure that is not a literal subalgebra of $H_q(S_n)$.
\end{enumerate}
We expand on (a)--(c) and the structural surprises uncovered along the
way.
\end{abstract}

\section{Setup}\label{sec:setup}

We retain the notation of the May-13 paper (\texttt{2026-05-13-multi-corner-dimension.tex}).
In particular: $H_q(S_n)$ is the Iwahori--Hecke algebra with generators
$T_1, \ldots, T_{n-1}$; $V_\lambda$ is the Specht module in the
Hoefsmit seminormal basis; $R'_j := (T_{j-1}+1)\cdots(T_1+1)$ and
\[
   B_{\ell+1}^{(\lambda)} V_\lambda \;:=\; R'_{\ell+1}R'_{\ell+2}\cdots R'_n V_\lambda
   \;\subseteq\; V_\lambda,
\]
the iterated-image subspace at the sharp index $\ell + 1$. The
decoration shape is $\lambda^{(\ge 2)} := (\lambda_1-1, \lambda_2-1, \ldots)$
(drop zero parts).

\medskip
The May-13 paper proves (modulo the technical ingredients there) that
in the stable regime,
\[
   \dim B_{\ell+1}^{(\lambda)} V_\lambda \;=\; f^{\lambda^{(\ge 2)}}.
\]
Since $|\lambda^{(\ge 2)}| = n - \ell$, the dimension matches that of
the Specht module $V_{\lambda^{(\ge 2)}}$, which is irreducible over
$H_q(S_{n-\ell})$ at generic $q$.

\begin{question}[Target]\label{q:target}
Is there a copy of $H_q(S_{n-\ell})$ inside $H_q(S_n)$ acting on
$B_{\ell+1}^{(\lambda)} V_\lambda$ such that the resulting module is
isomorphic to $V_{\lambda^{(\ge 2)}}$?
\end{question}

\noindent
The most natural candidates for such a copy are:
\begin{itemize}[leftmargin=2em,topsep=0.2em,itemsep=0.1em]
\item[(A)] the \emph{standard parabolic} embedding
$H_q(S_{n-\ell}) = \langle T_{\ell+1}, T_{\ell+2}, \ldots, T_{n-1}\rangle$
acting on positions $\{\ell + 1, \ldots, n\}$;
\item[(B)] the \emph{shifted parabolic} embedding
$H_q(S_{n-\ell}) = \langle T_2, T_3, \ldots, T_{n-\ell}\rangle$ acting on positions
$\{2, 3, \ldots, n - \ell + 1\}$ (or any other run);
\item[(C)] a copy generated by Jucys--Murphy elements, intertwiners, or
braid words.
\end{itemize}
Below we eliminate (A) and (B) outright by direct computation, then
discuss (C).

\section{Computational refutation}\label{sec:refute}

We work modulo $P = 100\,003$ at $q = 11$ in the Hoefsmit seminormal
matrix model (script \texttt{compute\_dim\_B.py}, May-13). The mod-$P$
rank equals the rank over $\mathbb{Q}(q)$ for any non-special $q$.

\subsection{The test case $\lambda = (3, 2, 1, 1, 1)$ at $n = 8$}

Here $\ell = 5$, $\dim V_\lambda = 64$, and
$\dim B_{6}^{(\lambda)} V_\lambda = 2 = f^{(2,1)}$ (where
$\lambda^{(\ge 2)} = (2, 1)$).

\begin{theorem}[Direct refutation]\label{thm:refutation}
For each $j \in \{1, 2, \ldots, 7\}$, computation gives:
\begin{center}
\renewcommand{\arraystretch}{1.15}
\begin{tabular}{c|c|l}
\toprule
$j$ & $T_j$ preserves $B$? & action / failure mode \\
\midrule
$1$ & \textbf{yes}  & scalar $-1$ \\
$2$ & no            & $\dim(B + T_2 B) = 4$ \\
$3$ & no            & $\dim(B + T_3 B) = 4$ \\
$4$ & no            & $\dim(B + T_4 B) = 4$ \\
$5$ & \textbf{yes}  & scalar $q$ \\
$6$ & no            & $\dim(B + T_6 B) = 4$ \\
$7$ & no            & $\dim(B + T_7 B) = 4$ \\
\bottomrule
\end{tabular}
\end{center}
Hence the only generators stabilising $B$ are $T_1$ and $T_5$, both
acting as scalars; the subalgebra they generate is the commutative
algebra spanned by $1, T_1, T_5, T_1 T_5$, isomorphic to the group
algebra $\mathbb{Z}[\zeta_2^2]$ of a Klein 4-group (after rescaling).
\end{theorem}

\begin{proof}[Proof (computational)]
Let $V_\lambda$ have its Hoefsmit seminormal basis $\{v_T : T \in
\mathrm{SYT}(\lambda)\}$, and build $T_j$ as a $64 \times 64$
mod-$P$ matrix. Build $\mathcal{B} := R'_6 R'_7 R'_8$ as a matrix
product. Apply Gaussian elimination to find a basis $W$ of
$\mathrm{Im}(\mathcal{B})$ (two columns).

For each $j$, form the $64 \times 2$ matrix $T_j W$ and compute the
rank of $[W | T_j W]$ (mod $P$). If this rank equals $2$, then
$T_j W \subseteq \mathrm{span}(W)$ (so $T_j$ preserves $B$); else it
strictly extends $B$. The table above is the result of this test for
each $j$.

For $j \in \{1, 5\}$, the resulting $2 \times 2$ action matrix on $W$
is a scalar $\pm 1$ resp.\ $q$, verified by direct comparison to $\pm I_2$
resp.\ $q I_2$. \qed
\end{proof}

\begin{corollary}[Standard parabolic embedding fails]\label{cor:A-fail}
The standard $H_q(S_{n-\ell}) = \langle T_{\ell+1}, \ldots, T_{n-1}\rangle
= \langle T_6, T_7\rangle$ does \emph{not} preserve $B$ (neither
$T_6$ nor $T_7$ does individually). Hence (A) is impossible.
\end{corollary}

\begin{corollary}[Shifted parabolic embeddings also fail]\label{cor:B-fail}
For any run $\langle T_a, T_{a+1}, \ldots, T_{a + n - \ell - 2}\rangle$
of $n - \ell - 1 = 2$ consecutive generators (i.e.\ $a \in \{1, 2, \ldots, 6\}$
with $a + 2 \le 8$), at least one generator in the run fails to preserve
$B$, since $T_2, T_3, T_4, T_6, T_7$ all fail to preserve $B$.

(Only $T_1, T_5$ work, and they are non-consecutive, so they don't
generate a Hecke algebra of type $A$.)
\end{corollary}

\subsection{No short-word extensions help}

\begin{proposition}[Short-word search]\label{prop:short-words}
For each ordered pair $(i, j)$ with $1 \le i, j \le 7$, the product
$T_i T_j$ preserves $B$ if and only if both $T_i$ and $T_j$ already
preserve $B$ individually (i.e.\ $i, j \in \{1, 5\}$). Similarly for
braid triples $T_i T_j T_i$ with $|i - j| = 1$: none of them preserve
$B$ (other than via individual factors, which doesn't apply here).
\end{proposition}

\begin{proof}[Proof (computational)]
Enumerate all $49$ products $T_i T_j$ and all $\sim 12$ braid triples
$T_i T_j T_i$ with $|i - j| \le 1$, $i \ne j$; for each, apply to
the $2$-dim basis $W$ of $B$ and check whether the result lies in $W$.
Output: only products with $i, j \in \{1, 5\}$ preserve $B$. Braid
triples (all involving at least one of $T_2, T_3, T_4, T_6, T_7$) all
fail. \qed
\end{proof}

\subsection{Multiple test cases confirm the pattern}

The same phenomenon — only $T_1, \ldots, T_k$ (acting as scalar $-1$,
for some small $k$) and $T_\ell$ (acting as scalar $q$) preserve $B$
— holds across multiple stable test cases:

\begin{center}
\renewcommand{\arraystretch}{1.15}
\begin{tabular}{lcccc}
\toprule
$\lambda$ & $n$ & $\ell$ & $T_j = -1$ stable & $T_j = q$ stable \\
\midrule
$(3, 2, 1, 1, 1)$       & 8  & 5 & $\{1\}$    & $\{5\}$ \\
$(3, 2, 1, 1, 1, 1)$    & 9  & 6 & $\{1, 2\}$ & $\{6\}$ \\
$(4, 2, 1, 1, 1, 1)$    & 10 & 6 & $\{1\}$    & $\{6\}$ \\
$(3, 3, 1, 1, 1, 1)$    & 10 & 6 & $\{1\}$    & $\{6\}$ \\
$(4, 3, 1, 1, 1)$       & 10 & 5 & $\emptyset$ & $\{5\}$ \\
\bottomrule
\end{tabular}
\end{center}

\noindent In every case, the stable generators are $\{T_1, \ldots, T_k\}
\cup \{T_\ell\}$ for some $k = k(\lambda) \in \{0, 1, 2\}$, and they all
act as scalars. The threshold $k$ depends on $\lambda$ in a non-obvious
way — it is not just $\ell - r - 1$ or $q - 2$; nor is it just the
rank-zero predictor.

\section{Why the naive iso fails}\label{sec:why}

\subsection{Comparison to the Littlewood--Richardson decomposition}

A natural place where one would expect to find $V_{\lambda^{(\ge 2)}}$
inside $V_\lambda$ is the
$V_{(1^\ell)} \boxtimes V_{\lambda^{(\ge 2)}}$-isotypic component of
$V_\lambda \!\downarrow_{H_q(S_\ell) \times H_q(S_{n-\ell})}$.

By Littlewood--Richardson, this isotypic component has multiplicity
$c^\lambda_{(1^\ell), \lambda^{(\ge 2)}}$. For our $\lambda$, the
skew shape $\lambda / (1^\ell)$ has one cell in each row of $\lambda$
(since each row contributes column~$1$ to $(1^\ell)$); this is a
vertical strip iff each row has $\le 1$ cell removed, which is
automatic, and the LR condition uniquely identifies
$\lambda^{(\ge 2)}$ as the partition obtained by deleting the strip.
Thus $c^\lambda_{(1^\ell), \lambda^{(\ge 2)}} = 1$, and the isotypic
component has dimension exactly $f^{\lambda^{(\ge 2)}}$ — matching
$\dim B$.

\begin{observation}[The fatal mismatch]\label{obs:mismatch}
The $V_{(1^\ell)} \boxtimes V_{\lambda^{(\ge 2)}}$-isotypic component
of $V_\lambda \!\downarrow_{H_q(S_\ell) \times H_q(S_{n-\ell})}$
\emph{is not} $B$. The isotypic component requires $T_j v = -v$ for
all $j \in \{1, \ldots, \ell - 1\}$ (the sign rep on the first
factor), but Theorem~\ref{thm:refutation} shows that $T_2, T_3, T_4$
\emph{do not even preserve} $B$.
\end{observation}

So there are two distinct $f^{\lambda^{(\ge 2)}}$-dimensional subspaces
of $V_\lambda$ at the relevant size:
\begin{itemize}[leftmargin=2em,topsep=0.2em,itemsep=0.1em]
\item the LR isotypic component $V_{(1^\ell)} \otimes V_{\lambda^{(\ge 2)}}
\subseteq V_\lambda$ (which carries an honest $H_q(S_{n-\ell})$-action via
$T_{\ell+1}, \ldots, T_{n-1}$, and an $H_q(S_\ell)$-action as the sign rep);
\item the iterated-image $B = B_{\ell+1}^{(\lambda)} V_\lambda$ (which
carries only the scalar action of $T_1, \ldots, T_k, T_\ell$).
\end{itemize}
These have the same dimension by an LR-style coincidence, but they
are genuinely different subspaces, and the second has a
fundamentally weaker Hecke-action structure than the first.

\subsection{Hoefsmit-support structure of $B$}

Inspecting the Hoefsmit-basis decomposition of vectors in $B$ for
$\lambda = (3, 2, 1, 1, 1)$ reveals that the support is broad and
mixed: a typical basis vector of $B$ has nonzero coefficients on $\sim
20$+ SYTs of $\lambda$, including SYTs in which the column-$1$ entries
are \emph{not} a prefix of $\{1, 2, \ldots, \ell\}$. For example,
\[
   T_0 = [(1,3,4)\,(2,5)\,(6)\,(7)\,(8)]
\]
appears in the support: entry $3$ is at $(1, 2)$ (not in column~$1$),
entry $5$ at $(2, 2)$, and entries $\{6, 7, 8\}$ in the bottom three
column-$1$ rows. Thus the column-$1$ spine carries $\{1, 2, 6, 7, 8\}$,
not $\{1, 2, 3, 4, 5\}$, demonstrating that $B$ is \emph{not}
spanned by SYTs with column~$1$ filled by the smallest entries.

This already implies that $B$ is not the $V_{(1^\ell)}$-isotypic
component (which \emph{is} spanned by SYTs with column~$1$ filled by
$\{1, \ldots, \ell\}$).

\section{Where might the iso live? The three resolutions}\label{sec:resolutions}

The numerical coincidence $\dim B = f^{\lambda^{(\ge 2)}}$ is real (May-13
closed-form theorem, verified across $9$ shapes in the stable regime).
What it does \emph{not} naturally upgrade to is a Hecke-equivariant iso
with $V_{\lambda^{(\ge 2)}}$ under any of the obvious embeddings (A) or
(B). The remaining possibilities:

\subsection*{(a) Exotic operators in $H_q(S_n)$}
There might exist $X_1, \ldots, X_{n-\ell-1} \in H_q(S_n)$, not of the
form $T_i$ or simple products, that preserve $B$ and satisfy the
Hecke relations of $H_q(S_{n-\ell})$. A candidate source is the family
of \emph{Jucys--Murphy elements}
$L_k = T_{k-1}T_{k-2}\cdots T_1 \cdot T_1 \cdots T_{k-1}$ together with
\emph{intertwining operators} $\Phi_{k}$ that exchange axial distances.
Neither was located in this session; an exhaustive search would
require enumeration of much longer words.

\subsection*{(b) Vector-space-only iso}
The chain bijection $\mathrm{Chain}(\lambda) \leftrightarrow
\mathrm{SYT}(\lambda^{(\ge 2)})$ (Section~3 of the May-13 paper) is a
genuine bijection between two combinatorial sets and induces an
explicit isomorphism of \emph{vector spaces}
$\mathbb{C}^{\mathrm{Chain}(\lambda)} \cong V_{\lambda^{(\ge 2)}}$.
This bijection, however, is purely combinatorial; it does not on its
own give a Hecke action on $B$.

The honest content of the May-13 paper is:
\[
   \dim B \;=\; \#\mathrm{Chain}(\lambda) \;=\; f^{\lambda^{(\ge 2)}},
\]
not an iso of representations.

\subsection*{(c) Subquotient / induced structure}
$B$ might be a sub-quotient of $V_\lambda$ as an $H_q(S_n)$-module
whose graded pieces include $V_{\lambda^{(\ge 2)}}$ (now an
$H_q(S_n)$-module via some restriction-of-scalars or pullback). Or
$B$ might be the image of an $H_q$-equivariant morphism from an
induced module $\mathrm{Ind}_{H_q(S_\ell) \times H_q(S_{n-\ell})}^{H_q(S_n)}
V_{(1^\ell)} \boxtimes V_{\lambda^{(\ge 2)}}$. Both possibilities are
consistent with the data but require representation-theoretic
machinery not developed here.

\section{Computational evidence: stability threshold (backup work)}\label{sec:stability}

While probing the iso theorem, we ran the suggested backup target
(\textsc{Prove.md}, Section ``Backup target''): the stability threshold.
The result is striking and \emph{contradicts} the May-13 paper's
empirical claim that ``$q \ge 2$ trailing $1$-rows suffices.''

\begin{theorem}[Stability threshold is shape-dependent and non-monotone]\label{thm:stab-shape-dep}
Let $\widehat\lambda$ be a multi-column base partition (no trailing
$1$-rows) and write $\lambda(\widehat\lambda, q) := (\widehat\lambda, 1^q)$.
The minimum $q$ for which $\dim B_{\ell+1}^{(\lambda)} V_\lambda = f^{\lambda^{(\ge 2)}}$
holds is \emph{not} a universal constant; it depends on $\widehat\lambda$.
Moreover, the matching set $\{q : \dim B = f^{\lambda^{(\ge 2)}}\}$ is
\emph{not} an up-set in $q$ in general.
\end{theorem}

\begin{proof}[Proof (computational table)]
Direct mod-$P$ computation gives the following table, where $\dim B$
is computed via \texttt{compute\_dim\_B.py} mod $P = 100\,003$ at
$q = 11$:

\begin{center}
\renewcommand{\arraystretch}{1.15}
\begin{tabular}{l|cccccc}
\toprule
$\widehat\lambda \;\backslash\; q$ & $0$ & $1$ & $2$ & $3$ & $4$ & $5$ \\
\midrule
$(3, 2)$, target $f^{(2,1)} = 2$
   & $2\,\checkmark$ & $2\,\checkmark$ & $2\,\checkmark$ & $2\,\checkmark$ & $2\,\checkmark$ & $2\,\checkmark$ \\
$(4, 2)$, target $f^{(3,1)} = 3$
   & $3\,\checkmark$ & $\mathbf{6}\,\times$ & $3\,\checkmark$ & $3\,\checkmark$ & $3\,\checkmark$ & --- \\
$(3, 3)$, target $f^{(2,2)} = 2$
   & $\mathbf{1}\,\times$ & $\mathbf{3}\,\times$ & $2\,\checkmark$ & $2\,\checkmark$ & $2\,\checkmark$ & --- \\
$(4, 3)$, target $f^{(3,2)} = 5$
   & $\mathbf{4}\,\times$ & $\mathbf{7}\,\times$ & $\mathbf{12}\,\times$ & $5\,\checkmark$ & --- & --- \\
$(5, 2)$, target $f^{(4,1)} = 4$
   & $\mathbf{5}\,\times$ & $\mathbf{8}\,\times$ & $\mathbf{12}\,\times$ & --- & --- & --- \\
$(3, 2, 2)$, target $f^{(2,1,1)} = 3$
   & $3\,\checkmark$ & $3\,\checkmark$ & $3\,\checkmark$ & --- & --- & --- \\
\bottomrule
\end{tabular}
\end{center}

\noindent
Entries marked $\checkmark$: $\dim B = f^{\widehat\lambda^{(\ge 2)}}$.
Entries marked $\times$: failure, with the actual $\dim B$ in bold.
\qed
\end{proof}

\begin{remark}[Non-monotonicity at $(4, 2)$]\label{rem:non-monotone}
The $(4, 2)$ family shows the matching set
$\{q : \dim B = f^{\lambda^{(\ge 2)}}\} = \{0, 2, 3, 4\}$ — \emph{not an up-set}:
$q = 0$ matches, $q = 1$ fails ($\dim = 6 \ne 3$), then $q \ge 2$ matches again. This refutes any monotone-in-$q$ stability hypothesis.
\end{remark}

\begin{remark}[Refining the May-13 threshold claim]\label{rem:thresh-refined}
The May-13 paper conjectured ``$q \ge 2$ suffices empirically.'' This is
\emph{wrong}: for $\widehat\lambda = (4, 3)$, the threshold is
$q \ge 3$ (the $q = 2$ case fails with $\dim = 12 \ne 5$); for
$\widehat\lambda = (5, 2)$, we have failure all the way through $q = 2$
and the threshold is $\ge 3$ (untested at $q \ge 3$).

The threshold should be conjectured shape-dependently. A first guess:
\[
   q_0(\widehat\lambda) \;\stackrel{?}{=}\; \ell(\widehat\lambda) + (\text{depth of recursion}),
\]
but this is heuristic and unverified.
\end{remark}

\begin{remark}[Failure dimensions also non-obvious]\label{rem:fail-dims}
At the non-matching cases, the actual $\dim B$ does not seem to match a
simple alternative formula. For example, $(4, 2)$ at $q = 1$ gives
$\dim = 6 = 2 \cdot 3 = 2 f^{(3, 1)}$, and $(4, 3)$ at $q = 2$ gives
$\dim = 12$ versus prediction $5$ (no obvious factor). The structure
of $B$ in the non-stable regime remains open.
\end{remark}

\begin{remark}[Strict stable regime exhibits the formula]\label{rem:stable-formula}
When $q$ exceeds the shape-dependent threshold $q_0(\widehat\lambda)$,
the formula $\dim B = f^{\lambda^{(\ge 2)}}$ holds and is robust — this
is consistent with the May-13 closed-form theorem in the stable regime,
just with a stricter (and correct) definition of stability.
\end{remark>

\section{What was learned}\label{sec:learned}

\begin{enumerate}[leftmargin=2em,topsep=0.3em,itemsep=0.2em]
\item The dim equality $\dim B = f^{\lambda^{(\ge 2)}}$ is real
(\textsc{May-13} paper, modulo Phase A generalisation) \emph{in the
strict stable regime}, but the strict regime is shape-dependent and
narrower than the May-13 paper estimated.

\item The naive Hecke-iso upgrade is \emph{false}: there is no run
$\langle T_a, T_{a+1}, \ldots\rangle$ of $n - \ell - 1$ consecutive
generators preserving $B$.

\item The May-13 strategy note (\textsc{Prove.md} Step 1) suggested that
``$B$ is fixed under $(T_{j-1}+1)$ for $j > \ell$'' — this is consistent
only in the weak sense that $T_\ell$ acts as scalar $q$ on $B$
($B \subseteq E_\ell^+$). Higher $T_j$ ($j > \ell$) do \emph{not}
preserve $B$ as a subspace.

\item The number $k(\lambda)$ such that $T_1, \ldots, T_k$ act as $-1$
on $B$ is shape-dependent and not simply expressible in $\ell, r, q$ alone.

\item The stability threshold $q_0(\widehat\lambda)$ is shape-dependent
and the matching set $\{q : \dim B = f^{\lambda^{(\ge 2)}}\}$ can be
non-monotone (e.g.\ $(4, 2)$: matches at $q = 0$, fails at $q = 1$,
matches at $q \ge 2$). This is a substantive refinement of the May-13
paper's claim.

\item The iso, if it exists, requires non-obvious operators (option (a))
or a subquotient / induced structure (option (c)). Locating those is the
next concrete research target.
\end{enumerate}

\section{Honest status}

The target theorem of \textsc{Prove.md} (Section ``The target theorem'')
is \emph{not} proved here; in its naive form (any obvious parabolic
embedding of $H_q(S_{n-\ell})$) it is \emph{refuted}. The dim-formula
program is robust on its own — the May-13 paper, with its closed-form
SYT count in the stable regime, is the strongest statement we currently
have. The rep-theoretic content remains conjectural and now sits as an
open question of structural type, not a numerical one.

\end{document}
